\documentclass[11pt]{article}
\usepackage{amsmath,amssymb,amsthm,hyperref}
\begin{document}
\title{Audit of Ambient-Independent Relative-Rank Counting}
\author{}
\date{November 24, 2026}
\maketitle

\paragraph{Status.}
This note audits a literature route. It proves neither accepted target
and imports no external theorem into Lean. The original model and its
definitions are unchanged.

\paragraph{The relevant real theorem.}
Basu and Parida, \emph{Bounds on the realizations of zero-nonzero patterns
and sign conditions of polynomials restricted to varieties and applications},
\href{https://arxiv.org/html/2411.11729v5}{arXiv:2411.11729v5}, Theorem 11,
bounds the maximum, over connected components of a semialgebraic set $W$,
of its minimum additive rank relative to a real algebraic family $\Delta$.
Writing $p=\dim\Delta$, $D=\deg\Delta$, and defining $W$ with $s$
polynomials of degree at most $\delta\geq2$, the lower bound is
\[
 \Omega\left(\frac{\log b_0(W)}
 {(p+1)(\log s+\log\delta+\log D+\log\log b_0(W))}\right).
 \tag{1}
\]
The hypothesis requires each irreducible component of $\Delta$ to have
its full real dimension at every real point. The theorem concerns real
sign conditions; it is not merely a zero-nonzero theorem. Its product-family
estimate, Proposition 5.1, retains an explicit factor $H$ inside the
logarithm through the bound
$(O(sH\delta D))^{3H(p+1)}$. Thus the ambient independence does not
remove the head-count logarithm.

\paragraph{A valid smooth lift and its actual scale.}
Here is an independent application check, including the geometric
hypothesis. Set $N=2^n$. Normalize any admissible affine denominator by
its positive value at the zero vertex, so
\[
 B_b(x)=1+\sum_{i=1}^n b_i x_i,
 \qquad A_a(x)=a_0+\sum_{i=1}^n a_i x_i.
\]
A free score bias can be absorbed into one numerator whenever there is
at least one head. Let $P(b)=\prod_{x\in\{0,1\}^n}B_b(x)$ and consider
the complex affine algebraic set with coordinates $(a,b,z,(y_x)_x)$
defined by
\[
 P(b)z=1,\qquad y_x B_b(x)=A_a(x)\quad\hbox{for every }x.
 \tag{2}
\]
Projection to $(a,b)$ is an algebraic isomorphism from (2) onto the
principal open subset $P\neq0$ of $\mathbb C^{2n+1}$. Its inverse is
$z=P^{-1}$ and $y_x=A_a(x)/B_b(x)$. Hence this set is irreducible,
smooth, real, and has dimension $p=2n+1$ at every real point.
All admissible heads lift into it; relaxing denominator positivity and
orientation only enlarges the available family.

A projective parametrization is obtained by multiplying the constant
coordinate and every displayed coordinate by $P$. Each resulting
coordinate polynomial has degree at most $N+1$. Pulling back generic
hyperplanes and applying the usual projective B\'ezout bound gives
\[
 \deg\Delta\leq(N+1)^{2n+1}.
 \tag{3}
\]
Base points cannot increase this upper bound on the degree of the image.
This degree estimate is an external standard algebraic-geometric step,
not a Lean-checked claim here.

In the ambient real vector space, take
\[
 W=\{(a,b,z,y):y_x\neq0\text{ for every }x\}.
\]
Its connected components are exactly the $2^N$ choices of strict signs
of $y$, times the unrestricted auxiliary coordinates. It has $s=N$
linear defining tests, which are also of degree at most $\delta=2$.
A sum of lifted heads has the desired score vector in its $y$
coordinates; the sums of the auxiliary coordinates are unrestricted in
$W$. Therefore a missing low-rank component really would give a hard
Boolean sign function. Standard finite bias strictification handles the
original convention allowing zero on false outputs.

Substituting (3) into (1) gives only
$\Omega(N/n^3)$, weaker than the already checked Warren baseline.
Even a hypothetical constant value of $D$ in this application leaves
$\log s=\log\log b_0(W)=\Theta(n)$ and $p=\Theta(n)$ in (1), producing
the scale $N/n^2$, not the accepted $N/n$ target. This is a limitation
of this direct application, not a lower-bound obstruction for other
image-side arguments.

\paragraph{A separate generic-shattering caution.}
Sontag's \emph{VC Dimension of Neural Networks}, Section 7.1,
\href{https://mathsci.kaist.ac.kr/~nipl/mas557/VCD_ANN_3.pdf}{defines a
generic-shattering dimension} using density of shattered input tuples
in the entire configuration space. Its bound by twice the number of
parameters plus one is not an ordinary VC bound. Failure of density
allows a nonempty open set of shattered tuples. In particular it gives
no conclusion for shattering the particular Boolean cube. For a family
of continuous scores, strict shattering itself persists in a small
neighborhood of a finite tuple, by taking a minimum over the finitely
many chosen sign witnesses. Thus a claim that all larger strictly
shattered sets have measure zero would require a different theorem.
\end{document}
